% BTP1 final report, reviewed in parts.
% The opening draft contains the title, abstract, introduction, Cassini
% background, and cited references.
% Paste this entire file into main.tex in the supplied Overleaf template.
% Compile with pdfLaTeX. No additional files are required for this part.
% The complete report will be 6--8 pages; this partial draft is shorter.
%
% Result provenance:
% Shift policies: visualizations/shift_policy_results.json, source 89bc23a.
% Estimate errors: visualizations/remaining_work/remaining_work_results.json,
%                 source bb5f6df.
% LCM study: research/lcm-scoring:visualizations/lcm_scoring/lcm_results.json,
%            source 2e8df9a.

\documentclass[conference]{IEEEtran}
\usepackage[T1]{fontenc}
\usepackage{cite}
\usepackage{amsmath,amssymb,amsfonts}
\usepackage{graphicx}
\usepackage{textcomp}
\usepackage{xcolor}

\begin{document}

\title{Evaluating Communication Shift Costs and Scoring Accuracy in a Cassini-Based Simulator}

% Author details are intentionally left as placeholders for two students.
\author{
\IEEEauthorblockN{Student One}
\IEEEauthorblockA{\textit{Department Name} \\
\textit{Institute Name} \\
Roll Number}
\and
\IEEEauthorblockN{Student Two}
\IEEEauthorblockA{\textit{Department Name} \\
\textit{Institute Name} \\
Roll Number}
}

\maketitle

\begin{abstract}
Training jobs can slow each other down when they send data over the same
network link at the same time. Cassini reduces this overlap by choosing where
jobs run and adjusting when they communicate. We build a simplified simulator
based on this approach and study when its timing adjustments are worthwhile.
Our shift-selection rule uses the work left in each job to check whether a
new delay is expected to shorten average completion time. On a 200-job
synthetic short-job workload, it reduces mean completion time by 9.30 percent
compared with selecting shifts from traffic overlap alone, averaged over
three seeds with the same jobs and placement. Further experiments test
inaccurate work estimates. A separate traffic check shows that examining only
a short part of a repeating pattern can give an overly favourable score.
The findings come from model-based experiments on synthetic or illustrative
workloads, rather than measurements from a deployed training cluster.
\end{abstract}

\begin{IEEEkeywords}
distributed training, network-aware scheduling, communication alignment,
discrete-event simulation, periodic traffic
\end{IEEEkeywords}

\section{Introduction}
\label{sec:introduction}

Training a model on several graphics processing units (GPUs) requires both
computation and communication. Workers calculate model updates and exchange
data so that training can continue. In a shared cluster, several jobs may
use the same network links. If their communication bursts occur together,
their combined demand can exceed the available bandwidth. Training can then
slow down even when sufficient compute resources have been assigned.

The timing of these bursts offers a scheduling opportunity. While one job
computes, another may be able to use the network. Coordinating them could
reduce competition for a link without increasing its capacity or changing
the model being trained. We chose communication scheduling as our problem
because it offers this scheduling-level way to improve resource sharing.

Cassini provides a concrete starting point for studying this idea. It extends
an existing job scheduler with information about repeating communication
patterns. It compares possible worker placements and calculates delays that
help jobs share network links at different times~\cite{cassini}. Implementing
these ideas gives us a testbed in which the waiting cost and traffic score
of each decision can be examined separately.

Our investigation focuses on the difference between a favourable traffic
pattern and a useful scheduling decision. Creating a new alignment requires
a job to wait. A long job may recover this cost through less competition
later; a short job may finish sooner if it keeps its current timing. The
decision also depends on the score used to predict traffic overlap. Examining
only the beginning of a repeating pattern saves calculation time, but can
miss overlap that occurs later.

These observations lead to two research questions:
\begin{enumerate}
    \item When does a timing shift save enough time to justify the wait it adds?
    \item How does a shorter traffic-evaluation interval affect score accuracy
    and the quality of the selected shifts?
\end{enumerate}

We address the first question with a rule that considers the work remaining
in each affected job before accepting a new shift. We compare it with the
starting simulator and two simpler controls. The score control accepts only
higher scores; the cooldown control limits how often shifts can change.
We also test errors in the remaining-work estimates. Every policy receives
identical jobs and fixed job-to-link placement. We address the second question with a separate
traffic checker that evaluates selected shifts over a complete repeating
pattern. This separates score validation from the simulator's completion-time
formula.

The completed work has three main contributions:
\begin{itemize}
    \item A reproducible fast simulator with explicit shift-wait accounting
    and controlled comparisons under identical placement.
    \item An evaluation of remaining-work shift decisions against score and
    cooldown controls, including tests with systematic estimate errors.
    \item An independent LCM scoring study that separates optimistic scores
    from losses in selected-shift quality and measures their calculation cost.
\end{itemize}

\section{Background: How Cassini Coordinates Jobs}
\label{sec:cassini-background}

\subsection{Computation, Communication, and Timing Shifts}

A training iteration can contain intervals of computation and intervals of
network communication. A \emph{traffic profile} describes the bandwidth a job
requests over time. When the iteration repeats, a similar profile may repeat
with it. Different jobs can have different iteration lengths and several
communication bursts within one iteration.

Consider a simplified example with two jobs, each computing for 60~ms and
communicating for 40~ms per iteration. Assume each job needs the full link
capacity while communicating. If they start together, both request
the network during the same final 40~ms. Delaying the second job by 40~ms
separates their communication intervals in this ideal repeating pattern.
This delay is a \emph{time shift}. The example explains the opportunity, but
also its cost: the second job must first wait for 40~ms.

\subsection{Placement and Compatibility}

\emph{Placement} determines which servers run a job's workers and therefore
which network links carry its traffic. Cassini evaluates candidate placements
provided by an underlying scheduler. It favours combinations whose traffic
patterns can be interleaved on shared links. Its compatibility score expresses
how well those patterns can fit together after suitable shifts~\cite{cassini}.

To calculate shifts, Cassini represents a repeating iteration as a circle.
Intervals on the circle describe its traffic demand; rotating the circle
changes the job's timing offset. For jobs with different periods, the common
representation uses their \emph{least common multiple} (LCM). For example,
40~ms and 60~ms periods repeat together every 120~ms. Larger combinations can
produce a much longer common period and a more expensive calculation.

\subsection{Coordinating Jobs Across Links}

A job can use several links and share each with different jobs. Separate
link-level decisions might therefore request different shifts for the same
job. Cassini represents these relationships with an \emph{affinity graph}.
Its two types of vertices represent jobs and links; an edge indicates that
a job uses a link. For an acyclic graph, its traversal method obtains one
consistent shift for each job. The published method rejects placement
candidates whose affinity graphs contain cycles. Server agents then apply
the selected delays to the next iteration and correct timing drift during
execution~\cite{cassini}.

\subsection{Connection to Our Study}

The published system already applies real delays to training iterations.
Our first study adds a completion-time test to the shift-selection rule in
our simplified simulator. The second study examines the 600~ms scoring limit
used in our implementation; this limit is not presented as a setting from
the original Cassini system.

Our workloads are synthetic or illustrative. The fast simulator advances
between job arrivals and completions and estimates slowdown from compatibility,
without executing individual training operations or network transfers. The
separate checker directly calculates requested bandwidth over time, but does
not validate the simulator's mapping from traffic scores to completion time.
We use these limits when interpreting the experimental results.

% Next reviewed part: Related Work and System Model.
% Later parts will add the model, methods, setup, results, and conclusion here.

\begin{thebibliography}{00}
\bibitem{cassini}
S.~Rajasekaran, M.~Ghobadi, and A.~Akella,
``CASSINI: Network-Aware Job Scheduling in Machine Learning Clusters,''
in \emph{Proc. 21st USENIX Symposium on Networked Systems Design and
Implementation (NSDI)}, 2024, pp.~1403--1420.
\end{thebibliography}

\end{document}
